\documentclass[12 pt, letterpaper]{hmcpset}
\usepackage{hw-preamble}
\setlist[enumerate, 1]{label = (\roman*)}

\name{Elliot Bobrow, Filiana Kostopoulou, Stella Shah, Jayden Thadani}
\class{MATH 189 --- Commutative algebra}
\assignment{HW 1}
\duedate{3rd September 2025}

\begin{document}

\begin{problem}[A \& M Ch. 1 \#1.]
	Let $x$ be a nilpotent element of a ring $A$. Show that $1 + x$ is a unit of $A$. Deduce that the sum of a nilpotent element and a unit is a unit.
\end{problem}

\begin{solution}
	Since $x$ is nilpotent, there exists an even $n$ with $x^{n} = 0$. Then $1 - x^{n} = 1$. Since $t = - 1$ is a root of the polynomial $1 - t^{n} \in \mathbb{Z}[t]$, the polynomial factors as $(1 + t)q(t)$. Under the unique homomorphism $\varphi : \mathbb{Z}[t] \to A$ with $t \mapsto x$ the factorisation becomes
	\[
		1 = 1 - x^{n} = (1 + x)\varphi(q(x))
	\]
	That is, $1 + x$ is a unit with multiplicative inverse $\varphi(q(x))$.

	In general, if $y$ is a unit in $A$, then $y^{n}$ is also a unit. Thus, $y^{n} - x^{n} = y^{n}$ is a unit (again, here $n$ is even, with $x^{n} = 0$). We factor $u^{n} - t^{n}$ as $(u + t) q(u, t)$ in $\mathbb{Z}[u, t]$ (again, we can factor because the polynomial vanishes for $u = - t$). Thus, under the unique homomorphism $\psi : \mathbb{Z}[u, t] \to A$ with $u \mapsto y$ and $t \mapsto x$ the factorisation becomes
	\[
		y^{n} = y^{n} - x^{n} = (y + x) \psi(q(u, t)).
	\]
	Thus, $y + x$ is a unit with multiplicative inverse $\psi(q(u, t)) y^{-n}$ .
\end{solution}

\pagebreak

\begin{problem}[A \& M Ch. 1 \#5.]
	Let $A$ be a ring and let $A[[x]]$ be the ring of formal power series $f = \sum_{n = 0}^{\infty} a_{n} x^{n}$ with coefficients in $A$. Show that
	\begin{enumerate}[label = (\roman*)]
		\item $f$ is a unit in $A[[x]]$ $\iff$ $a_{0}$ is a unit in $A$.

		\item If $f$ is nilpotent, then $a_{n}$ is nilpotent for all $n \ge 0$. Is the converse true?

		\item $f$ belongs to the Jacobson radical of $A[[x]]$ $\iff$ $a_{0}$ belongs to the Jacobson radical of $A$.

		\item The contraction of a maximal ideal $\mathfrak{m}$ of $A[[x]]$ is a maximal ideal of $A$, and $\mathfrak{m}$ is generated by $\mathfrak{m}^{\text{c}}$ and $x$.

		\item Every prime ideal of $A$ is the contraction of a prime ideal of $A[[x]]$.
	\end{enumerate}
\end{problem}

\begin{solution}
	\begin{enumerate}
		\item Suppose $f = \sum a_{n} x^{n}$ is a unit in $A[[x]]$. Then there is some $g = \sum b_{n} x^{n}$ such that $fg = 1$. That is, we have
			\[
				\sum_{n = 0}^{\infty} \sum_{k = 0}^{n} a_{k} b_{n - k} x^{n} = 1
			\]
			and in particular, $a_{0} b_{0} = 1$. Thus, $a_{0}$ is a unit in $A$.

			Suppose $a_{0}$ is a unit in $A$ and $f$ is any $\sum a_{n} x^{n} \in A[[x]]$. We want to show that there exists some $g = \sum b_{n} x^{n}$ so that $fg = 1$. In particular we want $b_{n}$ such that $a_{0} b_{0} = 1$ and $\sum_{k = 0}^{n} a_{k} b_{n - k} = 0$ for all $n > 0$. We can have $a_{0} b_{0} = 1$ since $a_{0}$ is a unit in $A$. We show that the rest of the coefficients are $0$ by induction. 

			Let $fg = \sum c_{n} ^{n}$. Note that for $b_{1} = - a_{1} b_{0}^{2}$ we get
			\[
				c_{1} = a_{0} b_{1} + a_{1} b_{0} = a_{0} b_{0} (-a_{1} b_{0}) + a_{1} b_{0} = 0.
			\]
			Suppose that for all $k \le n$ we have $b_{k}$ such that $c_{k} = 0$. Then we can get $b_{n + 1}$ such that $c_{n + 1} = 0$ by choosing
			\[
				b_{n + 1} = - a_{0}^{-1} \sum_{k = 1}^{n} a_{k} b_{n - k}.
			\]
			It follows then that
			\[
				c_{n + 1} = a_{0} b_{n + 1} + \sum_{k = 1}^{n} a_{k} b_{n - k} = 0
			\]
			as the sum of additive inverses. By induction, we can get such $b_{n}$ for each $n$, and thus, we get $g = \sum b_{n} x^{n}$ that is the multiplicative inverse of $f$.

		\item Suppose $f = \sum a_{n} x^{n}$ is nilpotent (say $f^{m} = 0$). We show by induction that each $a_{n}$ is nilpotent. 

			Since the $f^{m} \pmod{x} \equiv a_{0}^{m}$, we must have $a_{0}^{m} = 0$.

			Suppose that $a_{n}^{m_{n}} = 0$ for each $n \le N$.

		\item Write $\mathfrak{R}_{A}$ and $\mathfrak{R}_{A[[x]]}$ for the Jacobson radicals of $A$ and $A[[x]]$ respectively. We know that $f = \sum a_{n} x^{n} \in \mathfrak{R}_{A[[x]]}$ if and only if $1 - fg$ is a unit in $A[[x]]$ for all $g = \sum b_{n} x^{n} \in A[[x]]$. We know (by part (i)) that the latter is equivalent to $1 - a_{0} b_{0}$ being a unit in $A$ for all $b_{0} \in A$. This is equivalent to $a_{0} \in \mathfrak{R}_{A}$.

		\item Denote the unique ring homomorphism $A \to A[[x]]$ by $\varphi$. Suppose $\mathfrak{m} \subseteq A[[x]]$ is a maximal ideal.

		\item Again, denote the unique ring homomorphism $A \to A[[x]]$ by $\varphi$. Suppose $\mathfrak{p} \subseteq A$ is a prime ideal. Then its extension $\mathfrak{p}^{\text{e}}$ is all products $a f$ with $a \in \mathfrak{p}$ and $f \in A[[x]]$. 

			We claim $\mathfrak{p}^{\text{e}}$ is prime. Suppose $f g \in \mathfrak{p}^{\text{e}}$.
	\end{enumerate}
\end{solution}

\pagebreak

\begin{problem}[A \& M Ch. 1 \#7.]
	Let $A$ be a ring in which every element $x$ satisfies $x^{n} = x$ for some $n > 1$ (depending on $x$). Show that every prime ideal in $A$ is maximal.
\end{problem}

\begin{solution}
	We show that for each prime $\mathfrak{p} \subseteq A$, the quotient $A / \mathfrak{p}$ is a field. It follows that $\mathfrak{p}$ is maximal.

	Recall that since $\mathfrak{p}$ is prime, $A / \mathfrak{p}$ is an integral domain. Also recall that the canonical map $\pi : A \to A / \mathfrak{p}$ is a surjection --- for every $y \in A / \mathfrak{p}$ there is an $x \in A$ such that $\pi(x) = y$. Thus, each $y \in A / \mathfrak{p}$ satisfies
	\[
		y = \pi(x) = \pi(x^{n}) = y^{n}
	\]
	for some $n > 1$ (depending on $y$). We can rearrange this to say that for every $y$ there is some $n > 1$ such that $y(y^{n - 1} - 1) = 0$. Finally, since $A / \mathfrak{p}$ is an integral domain, this means that either  $y = 0$ or $y y^{n - 2} = 1$. That is, $y = 0$ or $y$ is a unit. Thus, $A / \mathfrak{p}$ is a field.
\end{solution}	

\pagebreak

\begin{problem}[A \& M Ch. 1 \#8.]
	Let $A$ be a ring $\neq 0$. Show that the set of prime ideals of $A$ has minimal elements with respect to inclusion.
\end{problem}

\begin{solution}
	We prove this by Zorn's lemma. Specifically, we show that the set of prime ideals is non-empty and every chain of prime ideals has a lower bound. Then by Zorn's lemma (applied to the dual poset of prime ideals ordered by inclusion), the set of prime ideals has minimal elements with respect to inclusion.

	Since $A \neq 0$, it contains a maximal ideal. Since maximal ideals are prime, it contains a prime ideal. That is, the set of prime ideals of $A$ is non-empty.

	Suppose $(\mathfrak{p}_{\alpha})_{\alpha \in \mathcal{I}}$ is a chain of prime ideals (so that either $\mathfrak{p}_{\alpha} \subseteq \mathfrak{p}_{\beta}$ or $\mathfrak{p}_{\beta} \subseteq \mathfrak{p}_{\alpha}$ given any $\alpha, \beta \in \mathcal{I}$). We claim $\bigcap_{\alpha \in \mathcal{I}} \mathfrak{p}_{\alpha}$ is a lower bound (in the poset of prime ideals ordered by inclusion). Clearly $\bigcap_{\alpha \in \mathcal{I}} \mathfrak{p}_{\alpha} \subseteq \mathfrak{p}_{\alpha}$ for each $\alpha \in \mathcal{I}$. It only remains to show that it is a prime ideal. Suppose $ab \in \bigcap_{\alpha \in \mathcal{I}} \mathfrak{p}_{\alpha}$. Then in each $\mathfrak{p}_{\alpha}$ we have $a \in \mathfrak{p}_{\alpha}$ or $b \in \mathfrak{p}_{\alpha}$. Since the $\mathfrak{p}_{\alpha}$ are ordered by inclusion, it follows that at least one of $a$ or $b$ is in all of the $\mathfrak{p}_{\alpha}$, and thus, in their intersection. That is, their intersection is prime.
\end{solution}

\pagebreak

\begin{problem}[A \& M Ch. 1 \#10.]
	Let $A$ be a ring, $\mathfrak{N}$ its nilradical. Show that the following are equivalent:
	\begin{enumerate}[label = (\roman*)]
		\item $A$ has exactly one prime ideal
		\item every element of $A$ is either a unit or nilpotent
		\item $A / \mathfrak{N}$ is a field.
	\end{enumerate}
\end{problem}

\begin{solution}
	We show (i) $\implies$ (ii) $\implies$ (iii) $\implies$ (i).

	Suppose $A$ has exactly one prime ideal $\mathfrak{p}$. Then $\mathfrak{N} = \mathfrak{p}$ (since it is the intersection of all prime ideals) and it is maximal since any larger maximal ideal would be a distinct prime. If $a \in A \setminus \mathfrak{N}$, then $(a, \mathfrak{N})$ (the ideal generated by $a$ and $\mathfrak{N}$) must be all of $A$ since $\mathfrak{N}$ is maximal. Thus, $ba + n = 1$ for some nilpotent $n \in \mathfrak{N}$ and some $b \in A$. Rearranging, we get $ba = 1 + n$. Then by Ch. 1 \#1, $ba$ is a unit. That is, $abc = 1$ for some $c \in A$. Thus, $a$ is a unit with inverse $bc$. Every other element $a \in \mathfrak{N} \subseteq A$ is a nilpotent, so we are done.

	Suppose every element of $A$ is either a unit or nilpotent. In particular, $\mathfrak{N}$ is the set of all nilpotents, and $A \setminus \mathfrak{N}$ the set of all units. Then $\mathfrak{N}$ is maximal since any larger ideal would include a unit, and thus, all of $A$. Thus, $A / \mathfrak{N}$ is a field.

	Suppsoe $A / \mathfrak{N}$ is a field. Then $\mathfrak{N}$ is maximal. However, $\mathfrak{N}$ the intersection of all prime ideals, and thus, contained in all of them. If there were more than one prime ideal, $\mathfrak{N}$ would be contained in a prime ideal distinct from itself, and not be maximal. Thus, there is only one prime ideal in $A$.
\end{solution}

\pagebreak

\begin{problem}[A \& M Ch. 1 \#11.]
	A ring $A$ is \emph{Boolean} if $x^{2} = x$ for all $x \in A$. In a Boolean ring $A$, show that
	\begin{enumerate}[label = (\roman*)]
		\item $2x = 0$ for all $x \in A$ 
		\item every prime ideal $\mathfrak{p}$ is maximal and $A / \mathfrak{p}$ is a field with two elements
		\item every finitely generated ideal in $A$ is principal.
	\end{enumerate}
\end{problem}

\begin{solution}
	\begin{enumerate}
		\item Since $A$ is Boolean, we have $(x + 1)^{2} = x + 1$. Thus, $x^{2} + 2x + 1 = x^{2} + 1$. Cancelling $x^{2} + 1$ on both sides, we get $2x = 0$.

		\item By Ch. 1 \#7, we get that in a Boolean ring, $A / \mathfrak{p}$ is a field where every element satisfies $y(y - 1) = 0$. Since $A / \mathfrak{p}$ is an integral domain, this forces $y = 1$ or $y = 0$. Since any field must have distinct $1$ and $0$, $A / \mathfrak{p}$ has exactly those two elements.

		\item Since $(x_{1}, \dots, x_{n}) = (x_{1}, \dots, x_{n - 1}) + (x_{n})$, by induction it suffices to show $(x, y) = (x) + (y)$ is actually a principal ideal $(z)$ for some $z \in A$. In particular, we claim $z = x + y - xy$ has $(z) = (x, y)$. Clearly, $(z) \subseteq (x, y)$ since $z \in (x, y)$. Conversely, it is easy to verify that $xz = x^{2} = x$ and similarly $yz = y$. Thus, $x, y \in (z)$ and $(x, y) \subseteq (z)$. By double containment, $(z) = (x, y)$.
	\end{enumerate}
\end{solution}

\pagebreak

\begin{problem}[A \& M Ch. 1 \#12.]
	A local ring contains no idempotent $\neq 0, 1$.
\end{problem}

\begin{solution}
	Suppose $A$ is a local ring with unique maximal ideal $\mathfrak{m}$. Let $x \in A$ be idempotent with $x \mapsto y$ under $A \to A / \mathfrak{m}$. Since homomorphisms preserve equations, we have $y^{2} = y$ in $A / \mathfrak{m}$. Rearranging, we get $y (y - 1) = 0$, and since $A / \mathfrak{m}$ is an field, either $y = 0$ or $y = 1$.

	If $y = 0$, then $x \in \mathfrak{m}$. Since $\mathfrak{m}$ is the only maximal ideal, we also have that $x \in \mathfrak{R}$ (the Jacobson radical of $A$). This implies $1 - x$ is a unit in $A$. Thus, the equation $x(1 - x) = 0$ is equivalent to $x = 0$.

	If $y = 1$, then $x = 1 + a$ for some $a \in \mathfrak{m}$. Since $x^{2} = x$, we must have  $1 + 2a + a^{2} = 1 + a$, and thus, $a^{2} + a = 0$. It follows that $-a \in \mathfrak{m}$ is idempotent. By the same argument as previously, $a = 0$, so $x = 1$.
\end{solution}

\pagebreak

\begin{problem}[A \& M Ch. 1 \#15. \emph{The prime spectrum of a ring}]
	Let $A$ be a ring and let $X$ be the set of all prime ideals of $A$. For each subset $E$ of $A$, let $V(E)$ denote the set of all prime ideals of $A$ which contain $E$. Prove that
	\begin{enumerate}[label = (\roman*)]
		\item if $\mathfrak{a}$ is the ideal generated by $E$, then $V(E) = V(\mathfrak{a}) = V(\sqrt{\mathfrak{a}})$. 
		\item $V(0) = X$, $V(1) = \text{\O}$.
		\item if $(E_{i})_{i \in I}$ is any family of subsets of $A$, then
			\[
				V \left ( \bigcup_{i \in I} E_{i} \right ) = \bigcap_{i \in I} V(E_{i})
			\]
		\item $V(\mathfrak{a} \cap \mathfrak{b}) = V(\mathfrak{a} \mathfrak{b}) = V(\mathfrak{a}) \cup V(\mathfrak{b})$ for any ideals $\mathfrak{a}, \mathfrak{b}$ of $A$.

			These results show that the sets $V(E)$ satisfy the axioms for closed sets in a topological space. The resulting topology is called the \emph{Zariski topology}. The topological space $X$ is called the \emph{prime spectrum} of $A$ and is written $\operatorname{Spec} A$.
	\end{enumerate}
\end{problem}

\begin{solution}
	\begin{enumerate}
		\item First we show $V(E) = V(\mathfrak{a})$ by double containment. Since $E \subseteq (E) = \mathfrak{a}$, any prime $\mathfrak{p}$ with $\mathfrak{a} \subseteq \mathfrak{p}$ has $E \subseteq \mathfrak{p}$. Thus, $V(\mathfrak{a}) \subseteq V(E)$. Since $\mathfrak{a}$, the ideal generated by $E$ is the smallest ideal containing $E$ (by definition), if we have $\mathfrak{p} \in V(E)$, we have $E \subseteq \mathfrak{a} \subseteq \mathfrak{p}$ and thus, $\mathfrak{p} \in V(\mathfrak{a})$. That is, $V(E) \subseteq V(\mathfrak{a})$.

			Now we show $V(\mathfrak{a}) = V(\sqrt{\mathfrak{a}})$ (again by double containment). Suppose $\mathfrak{p} \in V(\mathfrak{a})$. Thus, $\mathfrak{a} \subseteq \mathfrak{p}$. For any $x \in \sqrt{\mathfrak{a}}$, we have $x^{n} \in \mathfrak{a} \subseteq \mathfrak{p}$. But then, since $\mathfrak{p}$ is prime, we must have $x \in \mathfrak{p}$. Thus, $\sqrt{\mathfrak{a}} \subseteq \mathfrak{p}$ and $V(\mathfrak{a}) \subseteq V(\sqrt{\mathfrak{a}})$. Now suppose $\mathfrak{p} \in V(\sqrt{\mathfrak{a}})$. Since $\mathfrak{a} \subseteq \sqrt{\mathfrak{a}}$, we have $\mathfrak{a} \subseteq \sqrt{\mathfrak{a}} \subseteq \mathfrak{p}$ and $\mathfrak{p} \in V(\mathfrak{a})$. That is, $V(\sqrt{\mathfrak{a}}) \subseteq V(\mathfrak{a})$.

		\item Every ideal is an additive subgroup of $(A, +)$, and thus, contains $0$. That is every $\mathfrak{p} \in X$ has $0 \subseteq \mathfrak{p}$ and so, $\mathfrak{p} \in V(0)$. Thus, $V(0) = X$.

			Since $(1)$, the ideal generated by $1$ is the whole ring $A$ and primes $\mathfrak{p}$ must be proper subsets of $A$, there is no prime ideal containing $(1)$. That is $V(1) = \text{\O}$.

		\item Let $\mathfrak{a}_{i}$ be the ideals generated by each $E_{i}$. Then note that
			\[
				V\left ( \bigcup_{i \in I} E_{i} \right ) = V\left ( \sum_{i \in I} \mathfrak{a}_{i} \right )
			\]
			and
			\[
				\bigcap_{i \in I} V(E_{i}) = \bigcap_{i \in I} V(\mathfrak{a}_{i}).
			\]
			Thus, it suffices to prove
			\[
				V \left ( \sum_{i \in I} \mathfrak{a}_{i} \right ) = \bigcap_{i \in I} V(\mathfrak{a}_{i})
			\]
			which we will do by double containment.

			Suppose $\mathfrak{p} \in V(\sum_{i \in I} \mathfrak{a}_{i})$. For each $i$, we have $\mathfrak{a}_{i} \subseteq \sum_{i \in I} \mathfrak{a}_{i}$. Thus, $\mathfrak{a}_{i} \subseteq \sum_{i \in I} \mathfrak{a}_{i} \subseteq \mathfrak{p}$, and thus, $\mathfrak{p} \in V(\mathfrak{a}_{i})$ for each $i$. Thus, $\mathfrak{p} \in \bigcap_{i \in I} V(\mathfrak{a}_{i})$. 

			Suppose $\mathfrak{p} \in \bigcap_{i \in I} V(\mathfrak{a}_{i})$. Then each $\mathfrak{a}_{i} \subseteq \mathfrak{p}$. Since ideals are closed under addition, we must have $\sum_{i \in I} \mathfrak{a}_{i} \subseteq \mathfrak{p}$ and thus, $\mathfrak{p} \in V(\sum_{i \in I} \mathfrak{a}_{i})$.

		\item We show $V(\mathfrak{a} \cap \mathfrak{b}) = V(\mathfrak{a}) \cup V(\mathfrak{b})$ by ``triple containment''. That is, we show $V(\mathfrak{a} \cap \mathfrak{b}) \subseteq V(\mathfrak{a} \mathfrak{b}) \subseteq V(\mathfrak{a}) \cup V(\mathfrak{b}) \subseteq V(\mathfrak{a} \cap \mathfrak{b})$

			Suppose $\mathfrak{p} \in V(\mathfrak{a} \cap \mathfrak{b})$. Recall $\mathfrak{a} \mathfrak{b} \subseteq \mathfrak{a} \cap \mathfrak{b}$ (since all multiples of $a \in \mathfrak{a}, b \in \mathfrak{b}$ are in both $\mathfrak{a}$ and $\mathfrak{b}$). Thus $\mathfrak{a} \mathfrak{b} \subseteq \mathfrak{a} \cap \mathfrak{b} \subseteq \mathfrak{p}$ and $\mathfrak{p} \in V(\mathfrak{a} \mathfrak{b})$.

			Suppose $\mathfrak{p} \in V(\mathfrak{a} \mathfrak{b})$. If all $a \in \mathfrak{a}$ are also in $\mathfrak{p}$, then $\mathfrak{p} \in V(\mathfrak{a}) \subseteq V(\mathfrak{a}) \cup V(\mathfrak{b})$. Else, there exists some $a_{0} \in \mathfrak{a}$ with $a_{0} \notin \mathfrak{p}$. We have $a_{0} b \in \mathfrak{ab} \subseteq \mathfrak{p}$ for each $b \in \mathfrak{b}$. Since $\mathfrak{p}$ is prime, this forces that each $b \in \mathfrak{b}$ is also in $\mathfrak{p}$. Thus, $\mathfrak{b} \subseteq \mathfrak{p}$ and $\mathfrak{p} \in V(\mathfrak{b}) \subseteq V(\mathfrak{a}) \cup V(\mathfrak{b})$.

			Finally, suppose $\mathfrak{p} \in V(\mathfrak{a}) \cup V(\mathfrak{b})$.  Then  $\mathfrak{a} \cap \mathfrak{b} \subseteq \mathfrak{a} \subseteq \mathfrak{p}$ or $\mathfrak{a} \cap \mathfrak{b} \subseteq \mathfrak{b} \subseteq \mathfrak{p}$. In either case $\mathfrak{p} \in V(\mathfrak{a} \cap \mathfrak{b})$.
	\end{enumerate}
\end{solution}

\pagebreak

\begin{problem}[Lecture 1. \#1]
	For ring homomorphisms $f : A \to B$ and $g : B \to C$, show that $g \circ f : A \to C$ is a ring homomorphism.
\end{problem}

\begin{solution}
	Suppose $a_{1}, a_{2} \in A$ and $1_{A} \in A, 1_{B} \in B, 1_{C} \in C$ are the multiplicative identities. Then
	\begin{align*}
		g \circ f(a_{1} + a_{2}) & = g(f(a_{1} + a_{2})) \\
					 & = g(f(a_{1}) + f(a_{2})) \\
					 & = g(f(a_{1})) + g(f(a_{2})) \\
					 & = g \circ f(a_{1}) + g \circ f(a_{2})
	\end{align*}
	where the second equality follows because $f$ is a ring homomorphism and the third because $g$ is a homomorphism. Similarly,
	\begin{align*}
		g \circ f(a_{1} a_{2}) & = g(f(a_{1} a_{2})) \\
				       & = g(f(a_{1}) f(a_{2})) \\
				       & = g(f(a_{1})) g(f(a_{2})) \\
				       & = g \circ f(a_{1}) g \circ f(a_{2}).
	\end{align*}
	Finally,
	\begin{align*}
		g \circ f(1_{A}) & = g(f(1_{A})) \\
				 & = g(1_{B}) \\
				 & = 1_{C}.
	\end{align*}
\end{solution}

\pagebreak

\begin{problem}[Lecture 2. \#1]
	Find two different definitions of a subring.
\end{problem}

\begin{solution}
	Atiyah and Macdonald \cite[pg. 2]{atiyah-macdonald-1969} gives the following (incorrect) definition of a subring.
	\begin{definition}[subring]
		A subset $S$ of a ring $A$ is a \emph{subring} of $A$ if $S$ is closed under addition and multiplication and contains the identity element of $A$.
	\end{definition}

	Aluffi \cite[pg. 132]{aluffi-2009} gives the following more categorical definition.

	\begin{definition}[subring]
		A \emph{subring} $S$ of a ring $A$ is a ring whose underlying set is a subset of $A$ such that the inclusion function $S \xhookrightarrow{} A$ is a ring homomorphism.
	\end{definition}
\end{solution}

\bibliographystyle{alpha}
\bibliography{references}

\end{document}
